\begin{textsample}{NEG Dim 2 – 2021\_c6 – Score 1.00 – 2021\_c6/enade\_2021\_ci\_soc\_edcb3715cc.txt}  \label{ex:f2_neg_257}
CAPÍTULO 6 ANÁLISE TÉCNICA DA PROVA Este capítulo tem por objetivo apresentar o desempenho dos estudantes concluintes de Ciências Sociais ( Licenciatura/Bacharelado ) no Enade/2021.
Para isso, foram calculadas as estatísticas básicas da prova em seu todo ( seção 6.1.1 ), bem como as estatísticas relacionadas aos Componentes de Formação Geral ( seção 6.1.2 ) e de Conhecimento Específico da Área ( seção 6.1.3 ).
Dadas as suas características, foram analisadas, em separado, as questões objetivas ( seção 6.2 ) e as questões discursivas ( seção 6.3 ).
Tomando -se como base as duas Questões Discursivas do Componente de Formação Geral, nas seções 6.3.1.6 a 6.3.1.8, são apresentados os resultados e os comentários sobre a correção das respostas em relação à Língua Portuguesa.
Para as questões objetivas, foram disponibilizados os Índices de Facilidade e de Discriminação Ponto-Bisserial, também em separado, para os Componentes de Formação Geral ( seção 6.2.1 e 6.2.2 ) e de Conhecimento Específico ( seção 6.2.3 e 6.2.4 ).
De cada componente, uma das questões foi escolhida para exemplificar a análise gráfica, relacionando -se as alternativas escolhidas pelos estudantes ( inclusive o gabarito ) com o número de acertos no componente.
Nos Anexos IA e IB, apresenta -se a íntegra da análise gráfica para todas as questões objetivas, respectivamente à Licenciatura e ao Bacharelado.
Para cada uma das questões discursivas, os conteúdos dos tipos mais comuns de respostas dos estudantes são apresentados e comparados com o padrão de resposta esperado ( ver Anexos IXA e IXB com o padrão de respostas ).
Nas tabelas, constam as seguintes estatísticas das notas19 : média do desempenho na prova, erro padrão da média, desvio padrão, nota mínima, mediana e nota máxima para cada um de seus componentes.
Tais estatísticas contemplam o total de estudantes concluintes da Área de Ciências Sociais ( Licenciatura/Bacharelado ) inscritos e presentes à prova do Enade/2021, tendo em vista agregações, ou por Grandes Regiões e o país como um todo, segundo Habilitação e Modalidade de Ensino, ou por Categoria Administrativa, Organização Acadêmica da IES.
Em relação aos gráficos de distribuição de notas, o intervalo considerado foi de 10 unidades, aberto à esquerda e fechado à direita, com exceção do primeiro intervalo, [ 0 ; 10 ], fechado em ambos os extremos.
Para os gráficos de distribuição das notas das questões discursivas, foram consideradas mais duas categorias : questão em branco20 e nota zero. 19 Uma definição dessas estatísticas pode ser encontrada no Glossário. 20 Nesse grupo estão incluídas também as respostas classificadas como nulas ou desconsideradas. 6.1 ESTATÍSTICAS BÁSICAS DA PROVA Esta seção apresenta estatísticas selecionadas e histogramas da nota geral ( 6.1.1 ) e de cada componente : Formação Geral ( 6.1.2 ) e Conhecimento Específico ( 6.1.3 ).
São também apresentadas estatísticas selecionadas de subpopulações, caracterizadas por Grande Região, Categoria Administrativa, Organização Acadêmica e Modalidade de Ensino. 6.1.1 ESTATÍSTICAS BÁSICAS GERAIS Na Tabela 6.1, são apresentadas as Estatísticas Básicas da prova, por Grande Região segundo Habilitação e Modalidade de Ensino, dos estudantes concluintes de Ciências Sociais ( Licenciatura/Bacharelado ).
As análises serão feitas primeiramente para os cursos de Licenciatura e depois para os de Bacharelado.
Cada Habilitação é desagregada internamente em cursos Presenciais e cursos a Distância.
A Média das notas da prova, como um todo, dos alunos de Licenciatura ( nas seções seguintes serão analisados os Componentes de Formação Geral e de Conhecimento Específico ), foi 45,3, sendo que os estudantes da região Norte obtiveram a Média mais baixa ( 35,8 ), e os da região Sudeste obtiveram a Média mais alta ( 46,9 ).
As demais Médias foram : 42,5, na região Nordeste, 45,8, na região Sul e 44,8, na região Centro-Oeste.
O Desvio padrão para o Brasil, como um todo, foi 15,9, sendo o maior Desvio padrão encontrado na região Centro-Oeste ( 16,6 ), e o menor, na região Norte ( 9,4 ), indicando uma menor dispersão das notas nessa região.
A região que obteve a maior nota Máxima foi a Sul ( 89,7 ), ao passo que a região que atingiu a menor nota Máxima foi a Norte ( 55,2 ).
A Mediana do Brasil, como um todo, foi 45,9, sendo a maior Mediana obtida na região Sudeste ( 48,1 ), e a menor obtida na região Norte ( 35,8 ).
A nota Mínima do Brasil foi obtida nas regiões Nordeste, Sudeste e Sul, onde foi zero.
A maior nota Mínima, entre as cinco Grandes Regiões, foi 16,4, na região Norte e, a nota Mínima obtida na região Centro-Oeste foi 12,0.
Considerando -se as notas, como um todo dos alunos de Licenciatura segundo Grande Região, observa -se que existe diferença estatisticamente significativa ao nível de 95 \% entre a Média obtida na região Norte e as Médias obtidas nas demais regiões e, entre a Média obtida na região Nordeste e as Médias obtidas nas regiões Sudeste e Sul.
A Média das notas para os alunos de cursos Presenciais de Licenciatura, foi 48,1, sendo que os alunos da região Norte obtiveram a Média mais baixa ( 35,8 ), e os da região Sudeste obtiveram a Média mais alta ( 52,8 ).
As demais Médias foram : 42,8 na região Nordeste, 52,7 na região Sul e 44,9 na região Centro-Oeste.
O Desvio padrão das Licenciaturas Presenciais foi 15,5, sendo o menor obtido na região Norte 9,4 e o maior na região Centro-Oeste ( 16,6 ).
A maior Mediana foi a da região Sudeste ( 54,4 ) e a menor foi a da região Norte ( 35,8 ).
A maior nota Máxima foi 89,7, obtida na região Sul.
A nota Mínima foi zero obtida nas regiões Nordeste e Sudeste.
A maior nota Mínima, entre as cinco Grandes Regiões, foi 16,4, na região Norte e, as demais notas Mínima foram : 13,8 na região Sul e 12,0 no Centro-Oeste.
Considerando -se as notas segundo Grande Região, existe diferença estatisticamente significativa ao nível de 95 \% entre a Média da região Norte e as Médias das outras regiões.
A diferença entre a Médias das regiões Nordeste e Centro-Oeste e, entre as regiões Sudeste e Sul não são estatisticamente significativas.
Já a Média das notas da prova dos alunos de cursos a Distância de Licenciatura, 40,1, foi mais baixa do que a dos alunos de cursos Presenciais ( 48,1 ), sendo que os alunos da região Nordeste obtiveram a Média mais baixa ( 26,3 ), e os da região Centro-Oeste obtiveram a Média mais alta ( 41,7 ).
As demais Médias foram : 39,3 na região Sudeste e 41,1 na região Sul.
Para esta Habilitação não foram oferecidos cursos a Distância na região Norte.
O Desvio padrão foi 15,3, sendo o menor obtido na região Nordeste 8,5 e o maior na região Centro-Oeste ( 22,8 ).
A região que obteve maior nota Máxima foi a região Sul ( 84,0 ), ao passo que a região que atingiu a menor nota Máxima foi a Nordeste ( 40,3 ).
A Mediana do Brasil, como um todo foi 39,4, sendo a maior Mediana obtida na região Centro-Oeste ( 41,7 ), e a menor, obtida na região Nordeste ( 23,5 ).
A nota Mínima foi zero, obtida nas regiões Sudeste e Sul.
Considerando -se as notas segundo Grande Região, não existe diferença estatisticamente significativa ao nível de 95 \% entre a maior Média, obtida na região Centro-Oeste ( 41,7 ), e as médias das demais regiões.
Cumpre notar que ainda que as notas dos cursos Presenciais sejam comparáveis com as dos cursos a Distância dentro de uma mesma Habilitação, não é possível comparar as notas obtidas pelos alunos de Licenciatura e de Bacharelado, já que as provas são, em parte, diferentes.
A Média das notas da prova, como um todo dos alunos de Bacharelado ( nas seções seguintes serão analisados os Componentes de Formação Geral e de Conhecimento Específico ), foi 42,8, sendo que os estudantes da região Norte obtiveram a Média mais baixa ( 32,2 ), e os da região Sul obtiveram a Média mais alta ( 45,3 ).
As demais Médias foram : 39,7, na região Nordeste e, 44,3, nas regiões Sudeste e Centro-Oeste.
O Desvio padrão para o Brasil, como um todo, foi 13,4, sendo o maior Desvio padrão encontrado na região Centro-Oeste ( 13,6 ), e o menor, na região Norte ( 10,9 ), indicando uma dispersão um pouco menor das notas dessa última região.
A região que obteve a maior nota Máxima foi a Sul ( 82,2 ), ao passo que a região que atingiu a menor nota Máxima foi a Norte ( 59,4 ).
A Mediana do Brasil, como um todo, foi 42,8, sendo a maior Mediana obtida na região Sul ( 46,0 ), e a menor obtida na região Norte ( 31,8 ).
A nota Mínima do Brasil foi zero, obtida em quase todas as regiões, exceto a Centro-Oeste ( 6,8 ).
Considerando -se as notas de todo o Brasil dos alunos de Bacharelado, segundo Grande Região, observa -se que não existe diferença estatisticamente significativa ao nível de 95 \% entre as Médias, obtidas nas regiões Sudeste, Sul e Centro-Oeste.
Já as duas menores Médias, obtidas nas regiões Norte e Nordeste, têm diferença estatisticamente entre si e em relação às demais Médias.
Como não foram oferecidos cursos a Distância de Bacharelado, os resultados para os cursos Presenciais são os obtidos para o Brasil como um todo.
O comportamento das notas dos estudantes de todo o Brasil pode ser observado no Gráfico 6.1, que apresenta um histograma com a distribuição das mesmas, segundo a Habilitação.
Essa é uma distribuição unimodal para os alunos de cursos de Licenciatura e de Bacharelado, com ambas as modas no intervalo ( 40 ; 50 ].
Na Tabela 6.2, são apresentadas informações referentes à Média da nota final, desagregadas por Categoria Administrativa, por Organização Acadêmica, segundo a Habilitação.
A Média das notas da prova, como um todo, dos alunos de Licenciatura, foi 45,3 ( Tabela 6.1 ).
Em relação à Categoria Administrativa, os estudantes de Licenciatura das IES Públicas obtiveram Média mais alta ( 47,4 ) que a Média nacional, e os das IES Privadas obtiveram Média mais baixa ( 41,4 ).
Observa -se que existe diferença estatisticamente significativa entre as Médias das notas das IES Públicas e as das Privadas.
A diferença entre as Médias das regiões Sudeste e Norte ( 11,1 ), a maior e a menor Média, é superior à diferença entre a das IES Públicas e Privadas ( 6,0 ), caracterizando uma menor diversidade administrativa do que regional.
O Desvio padrão, igual para as IES Públicas e Privadas ( 15,6 ) foi um pouco inferior ao do Brasil, como um todo ( 15,9 ).
No tocante à Organização Acadêmica, as Universidades ( 46,1 ) obtiveram Média mais alta que a nacional ( 45,3 ).
A Média dos Centros Universitários e dos CEFET/IF foram menores que a nacional.
Constata -se que existe diferença estatisticamente significativa, ao nível de 95 \%, entre as Médias das Universidades e dos Centros Universitários.
Não foram oferecidos cursos de Licenciatura em Faculdades.
A Média das notas da prova, como um todo, dos alunos de Bacharelado, foi 42,8, mais baixa que a Média dos estudantes de Licenciatura ( Tabela 6.1 ).
Em relação à Categoria Administrativa, os estudantes de Bacharelado das IES Públicas obtiveram Média igual a Média nacional ( 42,8 ), e os das IES Privadas obtiveram Média muito próxima ( 42,7 ).
Observa -se que não existe diferença estatisticamente significativa entre as Médias das notas das IES Públicas e as das Privadas.
A diferença entre as Médias das regiões Sul e Norte ( 13,1 ), a maior e a menor Média, é bem superior à diferença entre a das IES Públicas e Privadas ( 0,1 ), caracterizando -se uma maior diversidade regional do que administrativa.
O Desvio padrão para as IES Públicas foi igual ao do Brasil, como um todo ( 13,4 ), não indicando diferença na dispersão das notas por Categoria Administrativa.
No tocante à Organização Acadêmica, as Universidades obtiveram Média ( 42,9 ) quase igual à nacional ( 42,8 ).
A Média dos Centros Universitários e a das Faculdades foram menores do que a nacional ( 34,7 e 40,6, respectivamente ).
Constata -se que não existe diferença estatisticamente significativa, ao nível de 95 \%, entre as Médias de todos os tipos de Organização Acadêmica.
Não foram oferecidos cursos de Bacharelado em CEFET/IFET. 6.1.2 ESTATÍSTICAS BÁSICAS NO COMPONENTE DE FORMAÇÃO GERAL Na Tabela 6.3, são apresentadas as Estatísticas Básicas em relação ao componente da prova que avalia a Formação Geral dos estudantes concluintes.
A nota de Formação Geral é obtida levando -se em conta os acertos nas oito questões objetivas para o curso de Ciências Sociais ( Licenciatura ).
Já para o curso de Ciências Sociais ( Bacharelado ), foram levadas em conta sete questões objetivas ( uma foi anulada pela CAA ) e a nota Média obtida nas Questões Discursivas 1 e 2, para ambas Habilitações.
A nota de cada questão discursiva é obtida levando -se em conta 80 \% da nota relativa ao conteúdo e 20 \% da nota relativa ao desempenho linguístico.
Os estudantes de Licenciatura, de todo o Brasil, obtiveram desempenho médio de 40,4.
Quanto à variabilidade, o Desvio padrão das notas dos estudantes de Licenciatura do Brasil foi 16,2.
A maior Média foi obtida na região Sudeste ( 42,5 ) e a menor, na região Norte ( 29,1 ).
As demais Médias foram : 38,0, na região Nordeste, 39,8, na região Sul e 40,1, na região Centro-Oeste.
Já o maior Desvio padrão foi obtido na região Sudeste ( 16,5 ), e o menor, na região Norte ( 13,6 ).
Os demais desvios foram : 15,0, na região Nordeste, 16,3, na região Sul e 15,9, na região Centro-Oeste.
A maior nota no Componente de Formação Geral da prova do Enade/2021 foi 94,0, obtida por, pelo menos, um estudante na região Sudeste.
A menor nota Máxima foi obtida na região Norte ( 57,7 ).
A Mediana do Brasil, como um todo, foi 40,4, sendo a menor Mediana encontrada na região Norte ( 30,6 ), e a maior, na região Sudeste ( 42,5 ).
A nota Mínima nessa parte foi zero em todas as regiões.
Considerando -se as notas, segundo Grande Região, observa -se que existe diferença estatisticamente significativa entre as Média das notas do Componente de Formação Geral da região Norte em relação as demais regiões e entre a região Sudeste e as regiões Nordeste e Sul.
Já os alunos de cursos Presenciais de Licenciatura obtiveram desempenho médio de 43,0.
Quanto à variabilidade, o Desvio padrão das notas do conjunto dos estudantes das Licenciaturas Presenciais foi 15,8.
A maior Média foi obtida na região Sudeste ( 48,2 ), e a menor, na região Norte ( 29,1 ).
As demais Médias foram : 38,2 na região Nordeste, 45,5 na região Sul e 40,2 na região Centro-Oeste.
O maior Desvio padrão foi obtido na região Centro-Oeste ( 15,9 ), e o menor, na região Norte ( 13,6 ).
Os demais desvios foram : 15,0 nas regiões Nordeste e Sudeste e, 15,2 na região Sul.
A nota Máxima no Componente de Formação Geral da prova de alunos de cursos Presenciais de Licenciatura foi 94,0, obtida por, pelo menos, um aluno na região Sudeste.
A menor nota Máxima foi obtida na região Norte ( 57,7 ).
A Mediana do conjunto de Licenciaturas Presenciais foi 43,1, sendo a menor Mediana encontrada na região Norte ( 30,6 ), e a maior, encontrada na região Sudeste ( 48,8 ).
A nota Mínima nesta parte foi zero em quase todas as regiões, exceto na região Sul ( 3,0 ).
Considerando -se as notas segundo Grande Região, observa -se que não existe diferença estatisticamente significativa entre as Média das notas das Licenciaturas Presenciais, no Componente de Formação Geral, das regiões Sudeste e Sul e, das regiões Nordeste e Centro-Oeste.
A Média das notas do Componente de Formação Geral dos alunos de Licenciatura a Distância foi 35,5, mais baixa do que a dos cursos Presenciais ( 43,0 ), o mesmo ocorrendo com a Mediana ( 35,5 ).
A nota Mínima foi zero, a mesma dos cursos Presenciais.
Para esta combinação de Modalidade e Habilitação, o Desvio padrão foi 15,7, e a nota Máxima foi 78,2, obtida por pelo menos um estudante na região Sul.
Não houve participação de estudantes de cursos a Distância de Licenciatura na região Norte.
Os estudantes de Bacharelado, de todo o Brasil, obtiveram desempenho médio de 41,2 ( Tabela 6.3 ).
Quanto à variabilidade, o Desvio padrão das notas dos estudantes de Bacharelado do Brasil, como um todo, foi 14,2.
A maior Média foi obtida na região Centro-Oeste ( 45,5 ), e a menor, na região Norte ( 36,8 ).
As demais Médias foram : 39,1, na região Nordeste ; 41,6, na região Sudeste ; e 42,1, na região Sul.
Já o maior Desvio padrão foi obtido na região Norte ( 15,2 ) e o menor, nas regiões Nordeste, Sudeste e Centro-Oeste ( 14,0 ).
Na região Sul o Desvio padrão foi 14,1.
A maior nota para os estudantes de Bacharelado no Componente de Formação Geral da prova do Enade/2021 foi 85,4, obtida por, pelo menos, um estudante na região Sudeste.
A menor nota Máxima foi obtida na região Norte ( 68,6 ).
A Mediana do Brasil, como um todo, foi 41,7, sendo a menor Mediana encontrada na região Norte ( 35,8 ), e a maior, na região Centro-Oeste ( 45,5 ).
A nota Mínima foi zero em quase todas as regiões, exceto na região Centro-Oeste ( 4,8 ).
Considerando -se as notas, segundo Grande Região, observa -se que existe diferença estatisticamente significativa entre as Médias das notas do Componente de Formação Geral obtidas nas regiões Norte e Nordeste e as das regiões Sudeste, Sul e Centro-Oeste, assim como entre as das regiões Sul e Centro-Oeste.
As diferenças entre as demais Médias não são estatisticamente significativas.
Como não foram oferecidos cursos a Distância de Bacharelado, os resultados para os cursos Presenciais são os obtidos para o Brasil como um todo.
O Gráfico 6.2 propicia a avaliação do desempenho dos estudantes do Componente de Formação Geral, a partir do histograma da distribuição das notas correspondentes.
A distribuição é unimodal para os alunos de cursos de Licenciatura e de Bacharelado, com ambas as modas no intervalo ( 40 ; 50 ].
Esse intervalo modal é o mesmo da distribuição de notas da prova, como um todo ( Gráfico 6.1 ).
Nota -se, ainda, que no Gráfico 6.2, as notas apresentam dispersão maior, quando comparadas a do Gráfico 6.1 ( distribuição das notas da prova ), confirmado pela comparação dos Desvios padrões que, em Formação Geral, aumenta de 15,9 ( Licenciatura ) para a nota da prova, como um todo ( Tabela 6.1 ), para 16,2 ( Licenciatura ) no Componente de Formação Geral ( Tabela 6.3 ).
No Bacharelado o Desvio padrão da nota da prova, como um todo ( Tabela 6.1 ), é de 13,4, enquanto que, para o Componente de Formação Geral ( Tabela 6.3 ), é de 14,2.
Na Tabela 6.4, são apresentadas as informações referentes ao desempenho dos concluintes do Componente de Formação Geral, em diferentes agregações : Categoria Administrativa, Organização Acadêmica e Habilitação.
Primeiramente, com respeito aos alunos de Licenciatura, observa -se que existe diferença estatisticamente significativa entre as Médias dos tipos de Categoria Administrativa : a maior Média foi obtida por estudantes de Licenciatura de IES Públicas ( 42,3 ) e a menor, pelos de IES Privadas ( 36,9 ).
Nota -se que, também, há diferença estatisticamente significativa entre as Médias das Universidades ( 41,1 ) e dos Centros Universitários ( 36,8 ), respectivamente a maior e a menor Média entre as Organizações Acadêmicas.
Não há diferença estatisticamente significativa entres as Médias dos CEFET/IF ( 39,8 ) e as dos demais tipos de Organização Acadêmica.
Como já mencionado, não foram oferecidos cursos de Licenciatura em Faculdades.
Entre os alunos de Bacharelado, observa -se que não existe diferença estatisticamente significativa entre as Médias dos tipos de Categoria Administrativa : 41,1 para estudantes de IES Públicas e 42,5 para estudantes de IES Privadas.
Nota -se que, também, não há diferença estatisticamente significativa entre as Médias dos quatro tipos de Organização Acadêmica : 41,2 para as Universidades, 40,4 para os Centros Universitários e 41,7 para as Faculdades.
Como mencionado anteriormente, não foram oferecidos cursos de Bacharelado em CEFET/IF. 6.1.3 ESTATÍSTICAS BÁSICAS DO COMPONENTE DE CONHECIMENTO ESPECÍFICO Na Tabela 6.5, são apresentadas as Estatísticas Básicas referentes ao Componente de Conhecimento Específico da Área de Ciências Sociais ( Licenciatura/Bacharelado ) por Grande Região, segundo Habilitação e Modalidade de Ensino.
A nota do componente de Conhecimentos Específicos leva em conta as notas da parte objetiva e a média das notas obtidas nas três questões discursivas, parte discursiva da Área.
A Média do desempenho dos estudantes de Licenciatura do Brasil, como um todo, foi 46,9.
A maior Média foi obtida na região Sudeste ( 48,4 ), e a menor, na região Norte ( 38,0 ).
As demais Médias foram : 44,0 na região Nordeste, 47,8 na região Sul e 46,4 na região Centro-Oeste.
Quanto à variabilidade das notas, o Desvio padrão do Brasil, como um todo para os alunos de Licenciatura, foi 17,7, sendo o maior Desvio padrão observado na região Centro-Oeste ( 18,5 ) e o menor, na região Norte ( 11,1 ).
Os demais desvios foram : 16,8 na região Nordeste, 17,7 na região Sudeste e 18,1 na região Sul.
A Mediana das notas dos estudantes de Licenciatura de todo o Brasil foi 47,8.
A maior Mediana ocorreu na região Sudeste ( 50,2 ), e a menor, na região Norte ( 36,9 ).
As demais Medianas foram : 43,3 na região Nordeste, 48,2 na região Sul e 46,4 na região Centro-Oeste.
A nota Máxima do Brasil, como um todo, foi 96,3, sendo obtida por, pelo menos, um estudante de Licenciatura na região Sul.
As demais notas Máximas foram : 60,2 na região Norte, 87,9 na região Nordeste, 92,2 na região Sudeste e 88,8 na região Centro-Oeste.
A nota Mínima foi zero nas regiões Nordeste, Sudeste e Sul.
As demais notas Mínimas foram : 19,3 na região Norte e 3,8 na região Centro-Oeste. região Nordeste, 87,9 na região Sudeste e 88,8 na região Centro-Oeste.
A nota Mínima foi zero nas regiões Nordeste, Sudeste e Sul.
As demais notas Mínimas foram : 19,3 na região Norte e 3,8 na região Centro-Oeste.
Considerando -se as notas das Licenciaturas do Componente de Conhecimento Específico segundo Grande Região, observa -se que existe diferença estatisticamente significativa entre a Média das notas das regiões Norte e Nordeste em relação a das regiões Sudeste e Sul, assim como entre a da região Norte e a da região Centro-Oeste.
Já para os alunos dos cursos Presenciais de Licenciatura, a Média do desempenho de Conhecimento Específico foi 49,8.
A maior Média foi obtida na região Sul ( 55,1 ), e a menor, na região Norte ( 38,0 ).
As demais Médias foram : 44,2 na região Nordeste, 54,4 na região Sudeste e 46,4 na região Centro-Oeste.
Quanto à variabilidade das notas, o Desvio padrão do conjunto das Licenciaturas Presenciais foi 17,5, sendo o maior Desvio padrão observado na região Centro-Oeste ( 18,5 ), e o menor, na região Norte ( 11,1 ).
Os demais desvios foram : 16,8 na região Nordeste, 16,3 na região Sudeste e 17,0 na região Sul.
A Mediana das notas dos estudantes de cursos Presenciais de Licenciatura de todo o Brasil foi 51,0.
A maior Mediana ocorreu na região Sul ( 56,2 ), e a menor, na região Norte ( 36,9 ).
As demais Medianas foram : 43,9 na região Nordeste, 55,3 na região Sudeste e 46,4 na região Centro-Oeste.
A nota Máxima das Licenciaturas Presenciais em Ciências Sociais foi 96,3, sendo obtida por, pelo menos, um aluno na região Sul.
As demais notas Máximas foram : 60,2 na região Norte, 87,9 na região Nordeste, 92,2 na região Sudeste e 88,8 na região Centro-Oeste.
A nota Mínima foi diferente de zero nas regiões Norte ( 19,3 ), Sul ( 6,8 ) e Centro-Oeste ( 3,8 ).
Considerando -se as Licenciaturas oferecidas presencialmente, observa -se que existe diferença estatisticamente significativa entre as duas maiores Médias das notas do Componente de Conhecimento Específico, obtidas nas regiões Sudeste e Sul, e as das demais regiões.
Para os alunos dos cursos a Distância de Licenciatura, a Média do desempenho de Conhecimento Específico foi 41,6.
A maior Média foi obtida na região Centro-Oeste ( 44,9 ), e a menor, na região Nordeste ( 27,0 ).
As demais Médias foram : 40,7 na região Sudeste e 42,8 na região Sul.
Quanto à variabilidade das notas, o Desvio padrão do conjunto das Licenciaturas a Distância foi 16,8, sendo o maior Desvio padrão observado na região Centro-Oeste ( 28,5 ), e o menor, na região Nordeste ( 8,6 ).
Os demais desvios foram : 16,4 na região Sudeste e 17,2 na região Sul.
A Mediana das notas dos estudantes de cursos a Distância de Licenciatura de todo o Brasil foi 40,3.
A maior Mediana ocorreu na região Centro-Oeste ( 44,9 ), e a menor, na região Nordeste ( 27,0 ).
As demais Medianas foram : 39,1 na região Sudeste e 42,4 na região Sul.
A nota Máxima das Licenciaturas a Distância em Ciências Sociais foi 95,4, sendo obtida por, pelo menos, um aluno na região Sul.
As demais notas Máximas foram : 40,3 na região Nordeste, 89,0 na região Sudeste e 65,0 na região Centro-Oeste.
A nota Mínima foi diferente de zero nas regiões Nordeste ( 15,5 ) e Centro-Oeste ( 24,7 ).
Considerando -se as Licenciaturas oferecidas a Distância, observa -se que não existe diferença estatisticamente significativa entre a maior Média das notas do Componente de Conhecimento Específico, obtida na região Centro-Oeste ( 44,9 ), e as das demais regiões, por outro lado, existe diferença estatisticamente significativa entre as Médias da região Nordeste e as das regiões Sudeste e Sul.
A Média do desempenho dos estudantes de Bacharelado do Brasil, como um todo, foi 43,3.
A maior Média foi obtida na região Sul ( 46,3 ), e a menor, na região Norte ( 30,7 ).
As demais Médias foram : 39,9, na região Nordeste ; 45,2 na região Sudeste ; e 43,9, na região Centro-Oeste.
Quanto à variabilidade das notas, o Desvio padrão do Brasil, como um todo para os alunos de Bacharelado, foi 15,3, sendo o maior Desvio padrão observado na região Centro-Oeste ( 15,6 ) e o menor, na região Norte ( 12,0 ).
Os demais desvios foram : 14,6, na região Nordeste, 15,2 na região Sudeste e 15,0, na região Sul.
A Mediana das notas dos estudantes de Bacharelado de todo o Brasil foi 43,1.
A maior Mediana ocorreu na região Sul ( 46,0 ), e a menor, na região Norte ( 29,8 ).
As demais Medianas foram : 39,1 na região Nordeste ; 45,1 na região Sudeste ; e 42,4 na região Centro-Oeste.
A nota Máxima do Brasil, como um todo, foi 93,7, sendo obtida por, pelo menos, um estudante de Bacharelado na região Sul.
As demais notas Máximas foram : 57,3, na região Norte ; 80,6, na região Nordeste ; 88,7, na região Sudeste e 88,9 na região Centro-Oeste.
A nota Mínima foi zero, obtida em todas as regiões.
Observa -se que existe diferença estatisticamente significativa entre as duas menores Médias, das notas do Componente de Conhecimento Específico, das regiões Norte e Nordeste e entre essas e as demais.
No entanto, não existe diferença estatisticamente significativa entre as três Médias mais altas, obtidas nas regiões Sudeste, Sul e Centro-Oeste.
Como não foram oferecidos cursos a Distância de Bacharelado, os resultados para os cursos Presenciais são os obtidos para o Brasil como um todo.
Assim como nos Gráficos 6.1 e 6.2, no Gráfico 6.3, apresentado a seguir, é possível ser feita uma avaliação do desempenho de concluintes de Ciências Sociais, segundo a Habilitação, em relação ao Componente de Conhecimento Específico, com um histograma da distribuição das notas correspondentes.
Esta também é uma distribuição unimodal para ambas as Habilitações, e o grupo modal é o ( 40 ; 50 ] para Bacharelado, o mesmo do grupo modal da prova, como um todo ( Gráfico 6.1 ) e ( 50 ; 60 ] para Licenciatura.
Na Tabela 6.6, é apresentada uma comparação dos resultados em relação a Categoria Administrativa, a Organização Acadêmica e segundo a Habilitação, agora levando -se em conta o desempenho de estudantes no Componente de Conhecimento Específico da prova.
No que se refere a Organização Acadêmica, para os alunos de Licenciatura, a maior Média foi a das Universidades ( 47,7 ), vindo a seguir a dos CEFET/IF ( 43,3 ) e, por fim, a dos Centros Universitários ( 43,1 ).
Como já comentado, não foram oferecidos cursos de Licenciatura em Faculdades.
O maior Desvio padrão, superior ao valor para o Brasil, como um todo, foi o dos CEFET/IF ( 18,0 ).
As Universidades obtiveram a maior nota Máxima ( 96,3 ).
Os Centros Universitários obtiveram nota Máxima 87,9, e os CEFET/IF, 80,0.
As Medianas foram : 48,9 nas Universidades, a maior, 42,9 nos Centros Universitários e 42,3 nos CEFET/IF.
A nota Mínima foi zero para quase todas as Organizações Acadêmicas, exceto para os CEFET/IF ( 19,3 ).
Observa -se que existe diferença estatisticamente significativa, ao nível de 95 \%, no Componente de Conhecimento Específico entre as notas obtidas nas Universidades e nos Centros Universitários.
Quanto a Categoria Administrativa, observa -se um comportamento semelhante àquele da parte de Formação Geral e à prova, como um todo, ou seja, existe diferença estatisticamente significativa entre as Médias das IES Públicas ( 49,1 ), a maior, e as das IES Privadas ( 42,9 ).
No que se refere a Organização Acadêmica, para os alunos de Bacharelado, a maior Média foi a das Universidades ( 43,5 ), vindo a seguir das Faculdades ( 40,2 ) e, por fim, a dos Centros Universitários ( 32,7 ).
O maior Desvio padrão, superior ao valor para o Brasil, como um todo, foi o das Universidades ( 15,4 ).
As Universidades também obtiveram a maior nota Máxima ( 93,7 ).
Os Centros Universitários obtiveram nota Máxima 40,5, a menor delas e, as Faculdades, 79,8.
As Medianas foram : 43,4 nas Universidades, a maior, 33,6 nos Centros Universitários, a menor delas e, 38,3 nas Faculdades.
A nota Mínima foi zero somente para as Universidades, as demais : 23,2 nos Centros Universitários e 15,4 nas Faculdades.
Observa -se que não existe diferença estatisticamente significativa, ao nível de 95 \%, no Componente de Conhecimento Específico entre as Médias de todos os tipos de Organização Acadêmica.
Quanto a Categoria Administrativa, observa -se um comportamento diferente àquele da parte de Formação Geral e à prova, como um todo, ou seja, não existe diferença estatisticamente significativa entre as Médias das IES Públicas ( 43,4 ) e as das IES Privadas ( 42,8 ). 6.2 ANÁLISE DAS QUESTÕES OBJETIVAS Esta seção apresenta estatísticas selecionadas e histogramas das Questões Objetivas de Formação Geral para Licenciatura ( 6.2.1 ) e para Bacharelado ( 6.2.2 ) e de Conhecimento estatisticamente significativa entre as Médias das IES Públicas ( 49,1 ), a maior, e as das IES Privadas ( 42,9 ).
No que se refere a Organização Acadêmica, para os alunos de Bacharelado, a maior Média foi a das Universidades ( 43,5 ), vindo a seguir das Faculdades ( 40,2 ) e, por fim, a dos Centros Universitários ( 32,7 ).
O maior Desvio padrão, superior ao valor para o Brasil, como um todo, foi o das Universidades ( 15,4 ).
As Universidades também obtiveram a maior nota Máxima ( 93,7 ).
Os Centros Universitários obtiveram nota Máxima 40,5, a menor delas e, as Faculdades, 79,8.
As Medianas foram : 43,4 nas Universidades, a maior, 33,6 nos Centros Universitários, a menor delas e, 38,3 nas Faculdades.
A nota Mínima foi zero somente para as Universidades, as demais : 23,2 nos Centros Universitários e 15,4 nas Faculdades.
Observa -se que não existe diferença estatisticamente significativa, ao nível de 95 \%, no Componente de Conhecimento Específico entre as Médias de todos os tipos de Organização Acadêmica.
Quanto a Categoria Administrativa, observa -se um comportamento diferente àquele da parte de Formação Geral e à prova, como um todo, ou seja, não existe diferença estatisticamente significativa entre as Médias das IES Públicas ( 43,4 ) e as das IES Privadas ( 42,8 ). 6.2 ANÁLISE DAS QUESTÕES OBJETIVAS Esta seção apresenta estatísticas selecionadas e histogramas das Questões Objetivas de Formação Geral para Licenciatura ( 6.2.1 ) e para Bacharelado ( 6.2.2 ) e de Conhecimento Específico para Licenciatura ( 6.2.3 ) e para Bacharelado ( 6.2.4 ).
São, também, apresentadas e comparadas as médias das subpopulações caracterizadas por Grande Região. 21 6.2.1 COMPONENTE DE FORMAÇÃO GERAL – CIÊNCIAS SOCIAIS ( LICENCIATURA ) Na Tabela 6.7a, são apresentadas as Estatísticas Básicas relativas às oito Questões Objetivas do Componente da prova que abrange a Formação Geral dos estudantes de Licenciatura, segundo a Modalidade de Ensino.
A Média do Brasil foi 46,0.
A menor Média foi encontrada na região Norte ( 38,1 ), e a maior, na região Centro-Oeste ( 48,2 ).
As demais Médias foram : 43,5 na região Nordeste, 48,0 na região Sudeste e 45,1 na região Sul.
O Desvio padrão do Brasil foi 19,8, sendo o maior Desvio padrão encontrado na região Sul ( 20,6 ), e o menor Desvio padrão encontrado foi na região Nordeste ( 18,4 ).
As Medianas do Brasil, como um todo, e de quatro regiões ( Nordeste, Sudeste, Sul e Centro-Oeste ) foram iguais a 50,0.
Na região Norte a Mediana foi igual a 37,5.
A nota Máxima 100,0 foi alcançada em quase todas as Grandes Regiões, exceto na região Norte ( 87,5 ).
As notas Mínimas ( 0,0 ) foram iguais para todas as regiões.
Para os alunos de cursos Presenciais, a Média foi 48,9.
A menor Média foi encontrada na região Norte ( 38,1 ), e a maior, na região Sudeste ( 53,5 ).
As demais Médias foram : 43,7 na região Nordeste, 52,0 na região Sul e 48,3 na região Centro-Oeste.
O Desvio padrão foi 19,4, sendo o maior Desvio padrão encontrado na região Centro-Oeste ( 19,9 ), e o menor Desvio padrão encontrado foi na região Sudeste ( 17,8 ).
As Medianas de quatro regiões ( Nordeste, Sudeste, Sul e Centro-Oeste ) foram iguais a 50,0.
Na região Norte a Mediana foi igual a 37,5.
A nota Máxima 100,0 foi alcançada em quase todas as Grandes Regiões, exceto na região Norte ( 87,5 ).
As notas Mínimas ( 0,0 ) foram iguais para todas as regiões.
Já para os alunos de cursos a Distância, a Média foi 40,5.
A menor Média foi encontrada na região Nordeste ( 30,0 ), e a maior, na região Sudeste ( 40,8 ).
As demais Médias foram : 40,4 na região Sul e 37,5 na região Centro-Oeste.
O Desvio padrão foi 19,4, sendo o maior Desvio padrão encontrado na região Sul ( 19,8 ), e o menor Desvio padrão encontrado foi na região Centro-Oeste ( 0,0 ).
As Medianas de três regiões ( Sudeste, Sul e Centro-Oeste ) foram iguais a 37,5.
Na região Nordeste a Mediana foi igual a 25,0.
A nota Máxima, 100,0, foi alcançada na região Sul e, a nota Mínima zero foi obtida nas regiões Sudeste e Sul. 21 As fórmulas utilizadas para o coeficiente de correlação ponto-bisserial e para o índice de facilidade estão apresentadas no capítulo 1.
Na Tabela 6.8a, são apresentados o Índice de Facilidade e o Índice de Discriminação ( Ponto-Bisserial ) para cada uma das Questões Objetivas do Componente de Formação Geral.
Quanto ao Índice de Facilidade, foram usadas as seguintes cores para diferenciar o nível de dificuldade da questão :  Azul para as questões classificadas com índice Muito fácil ( \&gt ; =0,86 ), verde para as questões classificadas com índice Fácil ( 0,61 a 0,85 ), amarelo para as questões classificadas com Médio ( 0,41 a 0,60 ), vermelho para as questões classificadas com Difícil ( 0,16 a 0,40 ) e roxo para as questões classificadas com Muito difícil ( \&lt ; =0,15 ).
Já quanto ao Índice de Discriminação, foram usadas as seguintes cores para qualificar a questão :  As questões classificadas com índice Fraco receberam a cor vermelha ( \&lt ; =0,19 ), as classificadas com Médio receberam a cor amarela ( 0,20 a 0,29 ), as classificadas com Bom receberam a cor verde ( 0,30 a 0,39 ) e as classificadas com Muito bom ( \&gt ; =0,40 ) receberam a cor azul.
As Questões Objetivas do Componente de Formação Geral, segundo o Índice de Facilidade, foram assim avaliadas : nenhuma questão foi anulada pela Comissão Assessora de Área.
Nenhuma teve o Índice de Facilidade classificado como Muito fácil, e duas questões foram tidas como Fácil, por terem índice de acertos no intervalo entre 0,61 e 0,85.
Três questões foram consideradas com índice de dificuldade Médio, situando -se no intervalo entre 0,41 e 0,60 do Índice de Facilidade, ou seja, houve entre 41,0 \% e 60,0 \% de acertos.
Três questões apresentaram acertos na faixa entre 0,16 e 0,40 classificadas como Difícil.
Nenhuma questão apresentou menos de 15 \% de acertos, razão pela qual não houve questão classificada como Muito difícil.
O Índice de Facilidade variou de 0,18 a 0,66, e o de Discriminação, de 0,26 a 0,49.
Quanto ao Índice de Discriminação, cinco questões foram tidas como Muito bom, duas questões foram tidas como Bom e uma classificada como Médio.
Na Tabela 6.9a, é apresentada a distribuição das questões, levando -se em conta, simultaneamente, a classificação dos dois índices.
As cinco questões que tiveram Índice de Discriminação Muito bom figuraram em dois níveis de dificuldade Fácil e Médio : duas classificadas na categoria Fácil ( Questões 6 e 7 ) do Índice de Facilidade, e três, na categoria Médio ( Questões 1, 3 e 5 ).
Em particular, a Questão 3 foi a que apresentou o maior poder discriminatório, com índice 0,49, e foi considerada Médio em termos de Facilidade, com uma proporção de 0,56 acertos.
O máximo de acertos foi alcançado pela Questão 7, com um Índice de Facilidade de 0,66.
O Índice de Discriminação foi Bom para as Questões 2 e 4.
No Gráfico 6.4a, para exemplificar, analisa -se o comportamento da Questão 3 de Formação Geral.
Trata -se de uma questão considerada com nível Médio de dificuldade, em relação à Facilidade, e a que obteve o maior Índice de Discriminação dessa parte da prova.
Neste gráfico, cada uma das cinco curvas representa o percentual de respostas em determinada alternativa da questão, em função do número de acertos dos estudantes nessa parte da prova ( Formação Geral/Múltipla Escolha ), antes de possíveis eliminações pelo critério Ponto-Bisserial.
A curva em verde corresponde à alternativa B, a alternativa correta para essa questão.
Observa -se que entre os estudantes com menor número de acertos, nessa parte do exame, a situação mais frequente foi a escolha da alternativa E ( em vermelho ), incorreta.
Por exemplo, entre os estudantes que acertaram duas respostas, 11,2 \% escolheram a alternativa A ( em azul ), 27,5 \% escolheram a alternativa B ( em verde, gabarito ), 6,6 \% escolheram a alternativa C ( em preto ), 9,2 \% escolheram a alternativa D ( em laranja ) e 45,3 \% escolheram a alternativa E ( em vermelho ).
Dentre os que acertaram duas respostas entre as questões de múltipla escolha de Formação Geral, 0,0 \% deixou a questão em branco e 0,2 \% marcou mais de uma alternativa, invalidando a questão.
À medida que o número de acertos aumenta, indicando desempenho melhor nessa parte da prova, aumenta, concomitantemente, a proporção de estudantes que selecionaram a alternativa correta B, atingindo 100 \% para os estudantes com oito acertos.
Essa análise permite verificar como a questão discriminou os grupos de desempenho, justificando -se o alto índice ( 0,49 ) obtido na questão.
Cumpre notar que não é possível inferir deste gráfico nem o Índice de Facilidade, que seria uma média da proporção ponderada pela quantidade de estudantes com cada uma das notas, nem o Índice de Discriminação Ponto-Bisserial, por razão equivalente.
No caso extremo, no qual a grande concentração dos acertos dos estudantes fosse abaixo de quatro, o Índice de Facilidade seria obrigatoriamente abaixo de 60,0 \% ( neste exemplo ).
Caso a concentração fosse em seis acertos ou mais, o índice seria obrigatoriamente acima de 90,0 \%.
Os gráficos relativos às demais questões de Formação Geral constam do Anexo IA. 6.2.2 COMPONENTE DE FORMAÇÃO GERAL – CIÊNCIAS SOCIAIS ( BACHARELADO ) A Tabela 6.7b, são apresentadas as Estatísticas Básicas relativas às oito Questões Objetivas válidas do Componente da prova que abrange a Formação Geral dos estudantes de Bacharelado.
A Média do Brasil foi 51,4.
A menor Média foi encontrada na região Norte ( 44,2 ), e a maior, na região Centro-Oeste ( 56,5 ).
As demais Médias foram : 48,3, na região Nordeste, 52,0, na região Sudeste e 53,7, na região Sul.
O Desvio padrão do Brasil foi 19,7, sendo o maior Desvio padrão encontrado nas regiões Nordeste e Sudeste ( 19,6 ), e o menor, na região Centro-Oeste ( 18,7 ).
Os demais desvios foram nas regiões Norte e Sul ( 19,5 ).
As Medianas do Brasil, como um todo, e de quatro regiões ( Norte, Nordeste, Sudeste e Sul ) foram iguais a 50,0.
A Mediana da região Centro-Oeste foi 62,5.
A nota Máxima 100,0 foi alcançada em quase todas as Grandes Regiões, exceto na região Norte ( 87,5 ).
As notas Mínimas ( 0,0 ) foram iguais para todas as regiões.
Não foram ofertados cursos a Distância de Bacharelado.
Na Tabela 6.8b, são apresentados o Índice de Facilidade e o Índice de Discriminação ( Ponto-Bisserial ) para cada uma das Questões Objetivas do Componente de Formação Geral.
A convenção das cores é a mesma da Tabela 6.8a.
As Questões Objetivas do Componente de Formação Geral, segundo o Índice de Facilidade, foram assim avaliadas : das oito questões válidas, nenhuma teve o Índice de Facilidade classificado como Muito fácil, e quatro questões foram tidas como Fácil, por terem índice de acertos entre 0,61 a 0,85.
Uma questão foi considerada com índice de dificuldade Médio, situando -se no intervalo entre 0,41 e 0,60 do Índice de Facilidade, ou seja, houve entre 41,0 \% e 60,0 \% de acertos, enquanto três questões foram classificadas na categoria Difícil, situando -se no intervalo entre 0,16 e 0,40.
Nenhuma questão apresentou menos de 15 \% de acertos, razão pela qual não houve questão classificada como Muito difícil.
O Índice de Facilidade variou de 0,19 a 0,71 e o de Discriminação, de 0,32 a 0,49.
Seis das oito questões válidas, quanto ao Índice de Discriminação, foram tidas como Muito bom e duas como Bom.
Na Tabela 6.8b, são apresentados o Índice de Facilidade e o Índice de Discriminação ( Ponto-Bisserial ) para cada uma das Questões Objetivas do Componente de Formação Geral.
A convenção das cores é a mesma da Tabela 6.8a.
As Questões Objetivas do Componente de Formação Geral, segundo o Índice de Facilidade, foram assim avaliadas : das oito questões válidas, nenhuma teve o Índice de Facilidade classificado como Muito fácil, e quatro questões foram tidas como Fácil, por terem índice de acertos entre 0,61 a 0,85.
Uma questão foi considerada com índice de dificuldade Médio, situando -se no intervalo entre 0,41 e 0,60 do Índice de Facilidade, ou seja, houve entre 41,0 \% e 60,0 \% de acertos, enquanto três questões foram classificadas na categoria Difícil, situando -se no intervalo entre 0,16 e 0,40.
Nenhuma questão apresentou menos de 15 \% de acertos, razão pela qual não houve questão classificada como Muito difícil.
O Índice de Facilidade variou de 0,19 a 0,71 e o de Discriminação, de 0,32 a 0,49.
Seis das oito questões válidas, quanto ao Índice de Discriminação, foram tidas como Muito bom e duas como Bom.
Na Tabela 6.9b, é apresentada a distribuição das questões, levando -se em conta, simultaneamente, a classificação dos dois índices.
As seis questões que tiveram Índice de Discriminação Muito bom figuraram entre três níveis de dificuldade Fácil, Médio e Difícil : quatro classificadas na categoria Fácil ( Questões 1, 3, 6 e 7 ) do Índice de Facilidade, uma na categoria Médio ( Questão 5 ) e, uma na categoria Difícil ( Questão 4 ).
Em particular, a Questão 3 foi a que apresentou o maior poder discriminatório, com índice 0,49, e foi considerada Fácil em termos de Facilidade, com uma proporção de 0,65 acertos.
O máximo de acertos foi alcançado pela Questão 7, com um Índice de Facilidade de 0,71.
No Gráfico 6.4b, para exemplificar, analisa -se o comportamento da Questão 3 de Formação Geral.
Trata -se de uma questão considerada Fácil, em relação à Facilidade e a que obteve o maior Índice de Discriminação dessa parte da prova.
Neste gráfico, cada uma das cinco curvas representa o percentual de respostas em determinada alternativa da questão, em função do número de acertos dos estudantes nessa parte da prova ( Formação Geral/Múltipla Escolha ), antes de possíveis eliminações pelo critério Ponto-Bisserial.
A curva em verde corresponde à alternativa B, a alternativa correta para essa questão.
Observa -se que, entre os estudantes com menor número de acertos, nessa parte do exame, a situação mais frequente foi a escolha da alternativa E ( em vermelho ), incorreta.
Por exemplo, entre os estudantes que acertaram duas questões, 17,4 \% escolheram a alternativa A ( em azul ), 29,4 \% escolheram a alternativa B ( em verde, gabarito ), 7,0 \% escolheram a alternativa C ( em preto ), 10,9 \% escolheram a alternativa D ( em laranja ), e 35,3 \%, a E ( em vermelho ).
Dentre os que acertaram duas respostas entre as questões de múltipla escolha de Formação Geral, nenhuma questão foi marcada com mais de uma alternativa ou em branco.
À medida que o número de acertos aumenta, indicando desempenho melhor nessa parte da prova, aumenta, concomitantemente, a proporção de estudantes que selecionaram a alternativa correta B, atingindo 100 \% para os estudantes com sete ou mais acertos.
Essa análise permite verificar como a questão discriminou os grupos de desempenho, justificando -se o alto índice ( 0,49 ) obtido na questão.
Cumpre notar que não é possível inferir deste gráfico nem o Índice de Facilidade, que seria uma média da proporção ponderada pela quantidade de estudantes com cada uma das notas, nem o Índice de Discriminação Ponto-Bisserial, por razão equivalente.
No caso extremo, no qual a grande concentração dos acertos dos estudantes fosse abaixo de quatro, o Índice de Facilidade seria obrigatoriamente abaixo de, aproximadamente, 60,0 \% ( neste exemplo ).
Caso a concentração fosse em seis acertos ou mais, o índice seria obrigatoriamente acima de 90,0 \%.
Os gráficos relativos às demais questões de Formação Geral constam do Anexo IB. 6.2.3 COMPONENTE DE CONHECIMENTO ESPECÍFICO – CIÊNCIAS SOCIAIS ( LICENCIATURA ) Na Tabela 6.10a, são apresentadas as Estatísticas Básicas em relação às Questões Objetivas do Componente de Conhecimento Específico da prova de Ciências Sociais por Grande Região.
Nessa parte da prova, como será discutido ainda nesta seção, nenhuma das Questões Objetivas do Componente de Conhecimentos Específicos foi anulada pela CAA e cinco questões não foram usadas no cômputo das notas por terem sido descartadas pelo critério do Ponto-Bisserial.
Assim, as notas foram calculadas com base em 22 das 27 questões objetivas válidas de Conhecimentos Específicos.
A Média do Brasil deste componente foi 47,6.
A menor Média foi observada na região Norte ( 39,8 ), e a maior, na região Sudeste ( 48,9 ).
O Desvio padrão de todo o Brasil foi 18,5, sendo o menor Desvio padrão encontrado na região Norte ( 12,2 ), e o maior, na região Centro-Oeste ( 19,1 ).
A Mediana de todo o Brasil foi 50,0, o mesmo valor da Mediana encontrada nas regiões Sudeste, Sul e Centro-Oeste.
Nas regiões Norte ( 38,7 ) e Nordeste ( 45,5 ), os valores foram menores.
A nota Máxima da prova foi 100,00, obtida na Questão Objetiva do Componente de Conhecimento Específico, por, pelo menos, um estudante da região Sul.
Nas demais regiões, a nota Máxima foi : 95,5, nas regiões Nordeste e Sudeste, 90,9 na região Centro-Oeste e, 63,6 na região Norte.
A nota Mínima foi zero, obtida por, pelo menos, um estudante nas regiões Nordeste, Sudeste e Sul nas Questões Objetivas do Componente de Conhecimento Específico.
A nota Mínima, na região Norte, foi 22,7 e, na região Centro-Oeste, foi 4,5.
Para os alunos de cursos Presenciais de Licenciatura, a Média deste componente foi 50,4, sendo que os alunos da região Norte obtiveram a Média mais baixa ( 39,8 ), e os da região Sul obtiveram a Média mais alta ( 55,7 ).
As demais Médias foram : 45,0 na região Nordeste, 54,8 na região Sudeste e 47,4 na região Centro-Oeste.
O Desvio padrão de todo o Brasil foi 18,5, sendo o menor obtido na região Norte 12,2 e o maior na região Centro-Oeste ( 19,1 ).
A maior Mediana foi a das regiões Sudeste e Sul ( 54,5 ) e a menor foi a da região Norte ( 38,7 ).
A maior nota Máxima foi 100,0, obtida na região Sul.
A nota Mínima foi zero obtida nas regiões Nordeste e Sudeste.
A maior nota Mínima, entre as cinco Grandes Regiões, foi 22,7, na região Norte.
As demais notas Mínima foram : 4,5 nas regiões Sul e Centro-Oeste.
Já para os alunos de cursos a Distância de Licenciatura, a Média deste componente foi 42,4, sendo que os alunos da região Nordeste obtiveram a Média mais baixa ( 29,1 ), e os da região Centro-Oeste obtiveram a Média mais alta ( 47,8 ).
As demais Médias foram : 41,4 na região Sudeste e 43,8 na região Sul.
O Desvio padrão de todo o Brasil foi 17,5, sendo o menor obtido na região Nordeste 9,1 e o maior na região Centro-Oeste ( 28,9 ).
A Mediana do Brasil foi 40,9, a mesma da região Sudeste.
A maior Mediana foi a da região Centro-Oeste ( 47,8 ) e a menor foi a da região Nordeste ( 29,6 ).
A maior nota Máxima foi 95,5, obtida na região Sul.
A nota Mínima foi zero obtida nas regiões Sudeste e Sul.
A nota Mínima foi zero, obtida nas regiões Sudeste e Sul.
A maior nota Mínima, entre as cinco Grandes Regiões, foi 27,3, na região Centro-Oeste.
A nota Mínima na região Nordeste foi 18,2.
Na Tabela 6.11a, são apresentados os Índices de Facilidade e de Discriminação ( Ponto-Bisserial ) das Questões Objetivas do Componente de Conhecimento Específico, para os estudantes de Ciências Sociais de Licenciatura.
Para facilitar a diferenciação das questões, serão usadas as mesmas cores da Tabela 6.8a e 6.8b para as diferentes classificações dos Índices de Facilidade e de Discriminação.
Nenhuma questão foi anulada pela CAA, sendo assim, as análises a seguir dizem respeito às 27 questões válidas.
A partir do Índice de Facilidade obtido, pode -se concluir que 40,7 \% das questões objetivas da prova foram consideradas Difícil : das 27 questões válidas, 11 foram classificadas como Difícil.
Cinco questões foram classificadas como Fácil, e outras nove consideradas como Médio.
Duas questões foram classificadas como Muito difícil e nenhuma questão foi classificada como Muito fácil.
Já quanto aos Índices de Discriminação das Questões Objetivas do Componente de Conhecimento Específico da prova, obtém -se como resultado a seguinte classificação : quatro das 27 questões válidas foram consideradas boas, enquanto 12 delas tiveram Índice de Discriminação Muito bom.
Assim, para 16 em 27 questões, os Índices de Discriminação foram Bom ou Muito bom.
Dentre as demais, seis delas foram classificadas como Médio, e outras cinco, como Fraco, sendo 11, por conseguinte, a quantidade de questões nos dois patamares mais baixos de Discriminação.
Constata -se, assim, que a prova – no que se refere ao Componente de Conhecimento Específico – possuía boa capacidade de discriminar entre aqueles que dominam ou não o conteúdo.
O Índice de Facilidade variou de 0,12 a 0,70 e o de Discriminação, de -0,07 a 0,53.
Na Tabela 6.12a, é apresentada a distribuição das questões levando -se em conta, simultaneamente, a classificação dos dois índices.
Dentre as questões que alcançaram os maiores Índices de Discriminação, todas foram classificadas com Índice de Discriminação Muito bom, e o maior índice ocorreu nas Questões 19 e 30.
O índice de ambas foi 0,53, e, quanto ao Índice de Facilidade, ambas também foram classificadas como Fácil, com 69,0 \% dos estudantes marcando a opção correta na Questão 19 e 68,0 \% dos estudantes marcando a opção correta na Questão 30.
As Questões de número 9 e 20 foram as mais difíceis entre as 27 questões específicas válidas, ambas com baixo Índice de Facilidade, apenas 12 \% de acertos.
Essas questões apresentaram poder discriminatório muito baixo, -0,07 para ambas, o que comprova terem sido as mais difíceis para os estudantes.
Destacam -se, também, as Questões 10, 12 e 25, com Índices de Facilidade, em termos percentuais, que correspondem a 36,0 \%, 27,0 \% e 16,0 \%, respectivamente, de estudantes que as responderam acertadamente.
E, sendo também, 0,03, 0,03 e 0,07, respectivamente, seus Índices de Discriminação.
Tais questões foram, portanto, pelo critério Ponto-Bisserial, consideradas inadequadas.
Por isso, as Questões 9, 10, 12, 20 e 25 foram eliminadas do cômputo da nota final.
A título de exemplo das análises do comportamento das questões objetivas, no Gráfico 6.5a, analisa -se a Questão 19 do Componente de Conhecimento Específico.
Essa questão foi considerada, pelas respostas dos estudantes avaliados na prova, como uma questão com nível de Facilidade Fácil, 0,69, ou seja, 69,0 \% dos estudantes assinalaram, acertadamente, a opção D correspondente ao gabarito.
Como já comentado, seu Índice de Discriminação foi igual a 0,53, classificado como Muito bom, o maior valor de Discriminação.
Neste gráfico, cada uma das cinco curvas representa o percentual de respostas em determinada alternativa da Questão 19, em função do número de acertos dos estudantes nessa parte da prova, antes de possíveis eliminações de questões pelo critério Ponto-Bisserial.
A alternativa correta D, representada no gráfico pela curva em laranja, foi escolhida em maiores proporções pelos estudantes com desempenho melhor nessa parte da prova.
Já as alternativas incorretas, também denominadas distratores, foram selecionadas, principalmente, por aqueles com notas mais baixas.
No caso de Ciências Sociais ( Licenciatura ), como se observa no eixo horizontal do Gráfico 6.5a, nenhum estudante acertou mais do que 24 questões dentre as 27 questões objetivas válidas, considerando -se o gabarito original.
A eliminação das cinco questões com Índice de Discriminação Fraco do cômputo da nota é que possibilitou a existência de nota Máxima igual a 100,0 por, pelo menos, um estudante da região Sul ( visto na Tabela 6.10a ).
Observa -se que a soma não é 100 \%, por causa das questões não respondidas ou com mais de uma opção marcada.
Aqueles com nota zero, na sua quase totalidade, marcaram a opção B ( linha verde ).
A proporção de estudantes que selecionaram a resposta correta D aumenta gradativamente, chegando a atingir 100 \% para 20, 22, 23 e 24 acertos, enquanto a proporção dos que escolheram alternativas incorretas decai, a partir de seis acertos, em função do número de acertos nessa parte da prova.
Os gráficos relativos às demais questões do Conhecimento Específico constam do Anexo IA. 6.2.4 COMPONENTE DE CONHECIMENTO ESPECÍFICO – CIÊNCIAS SOCIAIS ( BACHARELADO ) Na Tabela 6.10b, são apresentadas as Estatísticas Básicas em relação às Questões Objetivas do Componente de Conhecimento Específico da prova de Ciências Sociais por Grande Região.
Nessa parte da prova, nenhuma questão foi anulada pela Comissão Assessora da Área.
Além disso, como será discutido ainda nesta seção, seis Questões Objetivas do Componente de Conhecimentos Específicos não foram usadas no cômputo das notas por terem sido descartadas pelo critério do Ponto-Bisserial.
Assim, as notas foram calculadas com base em 21 das 27 questões objetivas válidas de Conhecimentos Específicos.
A Média do Brasil deste componente foi 42,0.
A menor Média foi observada na região Norte ( 29,4 ), e a maior, na região Sul ( 45,7 ).
O Desvio padrão de todo o Brasil foi 16,1, sendo o menor Desvio padrão encontrado na região Norte ( 12,2 ), e o maior, na região Centro-Oeste ( 16,8 ).
A Mediana de todo o Brasil foi 42,9, sendo o maior na região Sul ( 47,6 ) e a menor na região Norte ( 28,6 ).
A nota Máxima da prova foi 95,2, obtida nas Questões Objetivas do Componente de Conhecimento Específico, por, pelo menos, um estudante da região Sul.
A nota Mínima ( 0,0 ) foi a mesma em todas as regiões do Brasil.
Na Tabela 6.11b, são apresentados os Índices de Facilidade e de Discriminação ( Ponto-Bisserial ) das Questões Objetivas do Componente de Conhecimento Específico, para os estudantes de Ciências Sociais de Bacharelado.
Para facilitar a diferenciação das questões, serão usadas as mesmas cores da Tabela 6.8a, 6.8b e 6.11a para as diferentes classificações dos Índices de Facilidade e de Discriminação.
A partir do Índice de Facilidade obtido, pode -se concluir que 48,1 \% das questões objetivas da prova foram consideradas Difícil, 13 das 27 questões válidas, e quatro como Muito difícil.
Quatro questões foram classificadas como Fácil, e outras seis consideradas como Médio.
Nenhuma questão foi classificada como Muito fácil.
Já quanto aos Índices de Discriminação das Questões Objetivas do Componente de Conhecimento Específico da prova, obtém -se como resultado a seguinte classificação : 11 das 27 questões foram consideradas como Bom, enquanto seis delas tiveram Índice de Discriminação Muito bom.
Assim, para 17 em 27 questões, os Índices de Discriminação foram Bom ou Muito bom.
Dentre as demais, quatro delas foram classificadas como Médio, e seis, como Fraco, sendo 10, por conseguinte, a quantidade de questões nos dois patamares mais baixos de discriminação.
Constata -se, assim, que a prova – no que se refere ao Componente de Conhecimento Específico – possuía média capacidade de discriminar entre aqueles que dominam ou não o conteúdo.
O Índice de Facilidade variou de 0,09 a 0,79 e o de Discriminação, de -0,07 a 0,45.
Na Tabela 6.12b, é apresentada a distribuição das questões levando -se em conta, simultaneamente, a classificação dos dois índices.
Dentre as questões que alcançaram os maiores Índices de Discriminação, uma delas obteve o maior deles, a de número 19.
Seu índice foi 0,45, e, quanto ao Índice de Facilidade, foi classificada como Fácil, com 79 \% dos estudantes marcando a opção correta.
A Questão de número 29 foi a mais difícil entre as 27 questões específicas válidas, com baixo Índice de Facilidade, apenas 9 \% de acertos.
Essa questão apresentou poder discriminatório muito baixo, -0,07, o que comprova ter sido a mais difícil para os estudantes.
Destaca -se, também, as Questões 9 e 20, com Índice de Facilidade 0,12 e 0,11, o que, em termos percentuais, corresponde a 12,0 \% e 11,0 \%, respectivamente, de estudantes que responderam acertadamente.
Já seus Índice de Discriminação foram : 0,03, para a Questão 9 e, 0,00 para a Questão 20.
Tais questões foram, portanto, pelo critério Ponto-Bisserial, consideradas inadequadas.
Por isso, as Questões 9, 10, 12, 20, 29 e 30 foram eliminadas do cômputo da nota final.
A título de exemplo das análises do comportamento das questões objetivas, no Gráfico 6.5b, analisa -se a Questão 19 do Componente de Conhecimento Específico.
Essa questão foi considerada, pelas respostas dos estudantes avaliados na prova, como uma questão com nível de facilidade Fácil, 0,79, ou seja, 79,0 \% dos estudantes assinalaram, acertadamente, a opção D ( em laranja ), correspondente ao gabarito.
Como já comentado, seu Índice de Discriminação foi igual a 0,45, classificado como Muito bom, o maior valor de discriminação.
Neste gráfico, cada uma das cinco curvas representa o percentual de respostas em determinada alternativa da Questão 19, em função do número de acertos dos estudantes nessa parte da prova, antes de possíveis eliminações de questões pelo critério Ponto-Bisserial.
A alternativa correta D, representada no gráfico pela curva em laranja, foi escolhida em maiores proporções pelos estudantes com desempenho melhor nessa parte da prova.
Já as alternativas incorretas, também denominadas distratores, foram selecionadas, principalmente, por aqueles com notas mais baixas.
No caso de Ciências Sociais ( Bacharelado ), como se observa no eixo horizontal do Gráfico 6.5b, nenhum estudante acertou mais do que 22 questões dentre as 27 questões objetivas válidas, considerando -se o gabarito original.
Mesmo com a eliminação das seis questões com Índice de Discriminação Fraco do cômputo da nota, nenhum participante atingiu a nota Máxima, igual a 100,0, nessa parte da prova ( visto na Tabela 6.10b ).
Observa -se que a soma não é 100 \%, por causa das questões não respondidas ou com mais de uma opção marcada.
A proporção de estudantes que selecionaram a resposta correta D aumenta gradativamente, chegando a atingir 100 \% para 19 ou mais acertos, enquanto a

% matched lemmas: 
\end{textsample}
